\section{Arquitectura orientada a eventos}
\label{sec:eventos}
\textit{Responsables: Mariano (bus y mensajería, PLT-04) y Benjamín (catálogo de eventos, ARQ-01).}

\subsection{Catálogo y contratos}
El MVP usa 16 eventos con el formato \texttt{lms.<contexto>.<entidad>.<hecho>.v<N>}. Sus esquemas
(JSON Schema Draft-04), un ejemplo de cada uno y los contratos OpenAPI de los seis microservicios viven
en el repositorio \texttt{nubeduca-contratos}, que es la fuente de verdad. El \texttt{detail} de cada
evento tiene siempre la forma \texttt{\{"metadata": \{...\}, "data": \{...\}\}}; \texttt{metadata}
lleva \texttt{eventId}, \texttt{tenantId}, \texttt{occurredAt}, \texttt{aggregateType},
\texttt{aggregateId}, \texttt{aggregateVersion}, \texttt{actorId} y, en los eventos de un curso,
\texttt{courseId}. Los eventos llevan \textit{payload} mínimo: ids y lo que cambió, sin nombres ni
correos.

La CI de \texttt{nubeduca-contratos} valida en cada PR que cada evento del catálogo tenga esquema y
ejemplo, que cada ejemplo cumpla su esquema, y compara cada esquema contra \texttt{main} para marcar en
rojo los cambios incompatibles (PLT-02). Como el Learner Lab bloquea EventBridge Schema Registry
(sección~\ref{sec:limitaciones}), esta validación en CI lo reemplaza.

La Tabla~\ref{tab:eventos} resume el catálogo (\texttt{eventos/catalogo.yaml}). Todos los eventos
viajan por el bus \texttt{lms-<stage>}; \texttt{grade.posted} se publica además en el tópico SNS FIFO de
notas. Casi todos salen de una operación de la API; las excepciones son \texttt{tenant.created}, que lo
emite \texttt{poblar\_datos.py} al crear cada tenant, y \texttt{file.scanned}, que lo emite la Lambda de
antivirus cuando S3 avisa que llegó un archivo. Tres eventos no tienen consumidor todavía
(\texttt{tenant.created}, \texttt{course.archived} y \texttt{grade.changed}): se publican igual porque
son hechos del dominio y un servicio nuevo puede suscribirse sin cambiar al productor.

\begin{table}[htbp]
\centering\footnotesize
\begin{tabularx}{\textwidth}{l l >{\raggedright\arraybackslash}X l}
\toprule
Evento (\texttt{detail-type}) & Productor & Consumidores & Canal \\
\midrule
\texttt{lms.users.tenant.created.v1} & ms-usuarios & --- & bus \\
\texttt{lms.users.term.started.v1} & ms-usuarios & ms-cursos & bus \\
\texttt{lms.users.user.provisioned.v1} & ms-usuarios & ms-cursos, ms-notificaciones & bus \\
\midrule
\texttt{lms.courses.course.created.v1} & ms-cursos & ms-contenidos & bus \\
\texttt{lms.courses.course.published.v1} & ms-cursos & ms-notificaciones & bus \\
\texttt{lms.courses.course.archived.v1} & ms-cursos & --- & bus \\
\texttt{lms.courses.enrollment.created.v1} & ms-cursos & ms-contenidos, ms-evaluaciones, ms-calificaciones, ms-notificaciones & bus \\
\texttt{lms.courses.enrollment.dropped.v1} & ms-cursos & ms-contenidos, ms-evaluaciones, ms-calificaciones, ms-notificaciones & bus \\
\texttt{lms.courses.announcement.published.v1} & ms-cursos & ms-notificaciones & bus \\
\midrule
\texttt{lms.content.file.scanned.v1} & ms-contenidos & ms-evaluaciones & bus \\
\texttt{lms.content.module.published.v1} & ms-contenidos & ms-notificaciones & bus \\
\midrule
\texttt{lms.assessment.assignment.published.v1} & ms-evaluaciones & ms-calificaciones, ms-notificaciones & bus \\
\texttt{lms.assessment.submission.submitted.v1} & ms-evaluaciones & ms-notificaciones & bus \\
\texttt{lms.assessment.submission.graded.v1} & ms-evaluaciones & ms-calificaciones & bus \\
\midrule
\texttt{lms.grading.grade.changed.v1} & ms-calificaciones & --- & bus \\
\texttt{lms.grading.grade.posted.v1} & ms-calificaciones & ms-notificaciones & bus + SNS FIFO \\
\bottomrule
\end{tabularx}
\caption{Los 16 eventos del MVP, con su productor, sus consumidores y su canal
(\texttt{eventos/catalogo.yaml} de \texttt{nubeduca-contratos}).}
\label{tab:eventos}
\end{table}

\subsection{Bus y mensajería (PLT-04)}
El stack \texttt{infra/bus} del repositorio \texttt{nubeduca-backend} crea por stage:
\begin{itemize}
  \item El bus de EventBridge \texttt{lms-<stage>}, con un Archive de 7 días para reprocesar eventos.
        Su nombre queda en el parámetro SSM \texttt{/nubeduca/<stage>/bus-name}, que cada microservicio
        lee al desplegar.
  \item El tópico SNS FIFO \texttt{nubeduca-<stage>-notas.fifo} para \texttt{grade.posted}, con
        \texttt{MessageGroupId = <tenantId>:<courseId>} y \texttt{MessageDeduplicationId = eventId}, de
        modo que las notas de un curso llegan en orden.
  \item Los \textit{exports} de CloudFormation \texttt{nubeduca-<stage>-bus-name}, \texttt{-bus-arn} y
        \texttt{-notas-topic-arn}.
\end{itemize}
Cada consumidor declara en su propio stack una regla del bus, su cola SQS, una DLQ
(\texttt{maxReceiveCount: 5}) y una alarma de CloudWatch sobre la DLQ. Para eso hay una plantilla en
\texttt{plantillas/consumidor/}.

\subsection{Outbox transaccional con EventBridge Pipes (ADR-008)}
La Lambda escribe la entidad y el ítem \texttt{EVENTO\#<eventId>} en una misma
\texttt{TransactWriteItems}, después de validar el \texttt{detail} contra su esquema. Si el evento no
cumple el contrato, no se guarda nada. Un EventBridge Pipe lee el DynamoDB Stream de la tabla, deja
pasar solo los \texttt{INSERT} cuyo \texttt{sk} empieza con \texttt{EVENTO\#} y entrega el evento al
bus usando los atributos \texttt{source}, \texttt{detailType} y \texttt{detail} del ítem.

Durante GATE-0 se comprobó que \texttt{LabRole} permite crear Pipes con destino al bus. Con eso, el
equipo decidió en el ADR-008 usar Pipes por defecto (\texttt{outbox: pipe}) en staging y prod, en vez
de la Lambda relay: es lo que pedía el diseño original, elimina una Lambda por microservicio y no hace
falta transformar el evento. La Lambda relay sigue disponible como vuelta atrás con
\texttt{--param="outbox=lambda"}.

\subsection{Consumidores idempotentes}
La entrega es \textit{at-least-once}. El consumidor guarda \texttt{INBOX\#<eventId>} con escritura
condicional en la misma transacción que su cambio: si ya existía, descarta el evento como duplicado.
Además, si la \texttt{aggregateVersion} del evento es menor o igual a la ya vista (ítem
\texttt{AGG\#}), lo descarta por viejo. El consumidor es tolerante: ignora los campos que no conoce.

\subsection{Resultados y evidencias}
Verificado de punta a punta en AWS (GATE-0, 20-Nov-2026, cuenta de Dev 2, stage \texttt{dev}) con la
plantilla de servicio y \texttt{outbox=pipe}:
\begin{itemize}
  \item \texttt{POST /examples} guardó \texttt{EXAMPLE\#} y \texttt{EVENTO\#} en la misma transacción
        (201).
  \item El Pipe quedó en \texttt{RUNNING} y llevó el evento al bus \texttt{lms-dev}; la Lambda relay
        no se ejecutó (sin logs).
  \item La regla lo copió a la cola y el consumidor lo procesó (\texttt{INBOX\#<eventId>}, DLQ en 0).
  \item Al reenviar el mismo \texttt{eventId} con \texttt{put-events}, el consumidor lo descartó como
        \texttt{DUPLICADO} y siguió habiendo un solo \texttt{INBOX\#}.
\end{itemize}
Dev 1 repitió la prueba en staging con la relay (matrícula de ms-cursos que llega a
ms-notificaciones) y luego con Pipe. Finalmente, con los seis microservicios reales en staging, los
cinco Pipes \texttt{nubeduca-ms-*-staging-outbox} quedaron en \texttt{RUNNING}
(sección~\ref{sec:despliegue}).
